% TOP_PROBLEM: {"id":30007620,"problem_number":"LOCAL-30007620","title":"Home bias in cross-border portfolios","category_id":21,"category":{"id":21,"name":"economics","display_name":"Economics","description":"Open economic theory statements, published research agendas, and proposed scoped specifications; question provenance and historical priority are qualified per record.","slug":"economics","order_index":21,"created_at":"2026-10-08T00:00:00Z"},"status":"research_question","external_url":null,"legacy_ids":[],"published":false,"statement":"In the explicitly specified setting: How much of investors’ preference for domestic assets reflects rational hedging, taxes and investment barriers, information frictions or behavioral familiarity? The mathematical statement above fixes the target, assumptions and scope of this catalogue formulation.","economics":{"original_id":109,"field":"International finance and exchange rates","kind":"identification","formalization_basis":"adapted_research_question","source_status":"explicit_open","priority_status":"earliest_found","priority_note":"French–Poterba cites its own1990 article and Cooper–Kaplanis1986 holdings work. The verified question concerns equities; first priority for all cross-border holdings is not established.","earliest_verified_reference":{"authors":["Kenneth R. French","James M. Poterba"],"title":"Investor Diversification and International Equity Markets","year":1991,"url":"https://economics.mit.edu/sites/default/files/publications/2006858.pdf","doi":null,"locator":"Opening, p.222; Section III, pp.224–225","evidence_summary":"Explicitly asks why investors overweight domestic equity."},"current_reference":{"authors":["Bruno Pellegrino","Enrico Spolaore","Romain Wacziarg"],"title":"Barriers to Global Capital Allocation","year":2025,"url":"https://www.anderson.ucla.edu/faculty_pages/romain.wacziarg/downloads/2025_GCA.pdf","doi":"10.1093/qje/qjaf031","locator":"Abstract; Section II, manuscript p.14; Section III, pp.20–23, footnotes 12 and 15","evidence_summary":"Demand microfoundations remain nonunique; proposes information-linkage research."},"formalization_note":"The original verified question concerns equity home bias. Extending it to asset classes and causal mechanism shares is a proposed scope, with known rational hedging explanations allowed."},"statusQualification":"Published question or research agenda verified at the cited locator. The mathematical specification is adapted for this catalogue and includes additional assumptions; source publication does not establish first-ever priority or certify this exact model remains unsolved. The original verified question concerns equity home bias. Extending it to asset classes and causal mechanism shares is a proposed scope, with known rational hedging explanations allowed.","statusReviewedAt":"2026-10-08"}
% FIRST_POSTED: 2026-10-08
% ENTRY_KIND: statement_only
% !TeX program = lualatex
\documentclass[11pt]{article}
\usepackage{fontspec}
\usepackage[a4paper,margin=25mm]{geometry}
\usepackage{amsmath,amssymb,mathtools}
\usepackage{xurl}
\usepackage[hidelinks]{hyperref}
\newcommand{\sref}[2]{\hyperlink{source:#1:#2}{[S#2]}}
\setlength{\emergencystretch}{3em}
\begin{document}
% problemId:30007620
\section*{Home bias in cross-border portfolios}\label{30007620}
\subsection{Definitions and mathematical statement}
Fix investors from country \(c\), asset class \(j\) and date \(t\). Let \(w_{cjt}^{\mathrm{foreign}}\) be the observed share allocated to foreign claims and \(m_{cjt}^{\mathrm{foreign}}\) the foreign share of a prespecified investable world-market portfolio. For \(m_{cjt}^{\mathrm{foreign}}>0\), define
\[H_{cjt}=1-w_{cjt}^{\mathrm{foreign}}/m_{cjt}^{\mathrm{foreign}}.\]
This index compares holdings with a benchmark, not automatically with an investor's utility-maximizing allocation. A structural demand model includes domestic-income and exchange-rate hedging needs, taxes and transaction costs, legal restrictions, information-acquisition costs and familiarity-related belief distortions. The task is to identify the causal contribution of each mechanism to \(H_{cjt}\), using quantities, returns, investor characteristics and plausibly exogenous changes in those mechanisms. Define each contribution as the difference under a specified mechanism-removal intervention with other equilibrium responses retained. Address interactions and observational equivalence: different microfoundations may generate the same demand equation. Restate securities by ultimate issuer nationality when tax-haven domicile distorts geographic exposure, and distinguish domestic multinational exposure from directly foreign holdings. Report bounds where beliefs or hedging motives are unobserved. The classic equity phenomenon does not establish that every cross-border asset displays the same bias or that all deviations from world weights are irrational.

\par\textbf{Formulation and scope.} The original verified question concerns equity home bias. Extending it to asset classes and causal mechanism shares is a proposed scope, with known rational hedging explanations allowed.

\par\textbf{Open-question provenance.} An explicit question or research agenda was verified at the source locator. This does not by itself certify that every later version remains unresolved \sref{30007620}{1}.

\par\textbf{Historical priority.} French–Poterba cites its own1990 article and Cooper–Kaplanis1986 holdings work. The verified question concerns equities; first priority for all cross-border holdings is not established.

\subsection{Short English statement}
In the explicitly specified setting: How much of investors’ preference for domestic assets reflects rational hedging, taxes and investment barriers, information frictions or behavioral familiarity? The mathematical statement above fixes the target, assumptions and scope of this catalogue formulation.

\subsection{Sources}
\begin{itemize}\raggedright
\item[\textnormal{[S1]}]\hypertarget{source:30007620:1}{} Kenneth R. French, James M. Poterba. 1991. Investor Diversification and International Equity Markets. Opening, p.222; Section III, pp.224–225.
\url{https://economics.mit.edu/sites/default/files/publications/2006858.pdf}.
{\small\textit{Earliest verified question/agenda reference; historical priority qualified below. Explicitly asks why investors overweight domestic equity.}}

\item[\textnormal{[S2]}]\hypertarget{source:30007620:2}{} Bruno Pellegrino, Enrico Spolaore, Romain Wacziarg. 2025. Barriers to Global Capital Allocation. Abstract; Section II, manuscript p.14; Section III, pp.20–23, footnotes 12 and 15.
\url{https://www.anderson.ucla.edu/faculty_pages/romain.wacziarg/downloads/2025_GCA.pdf}.
DOI: 10.1093/qje/qjaf031.
{\small\textit{Supporting or current literature; see evidence qualification. Demand microfoundations remain nonunique; proposes information-linkage research.}}
\end{itemize}
\end{document}
